\clearpage
\appendix
\section{十分钟复现清单}

本清单以 Windows 源码模式为例，其他系统的差异只在虚拟环境激活命令。目标是在无 API Key、无 GPU 的条件下看到完整主路径。依赖下载时间受网络影响，表中的十分钟指常用网络和本机缓存条件下的操作预算，不把它作为所有机器的硬保证。追求最短现场启动时，优先使用已解压的 Windows 目录版。

\begin{longtable}{@{}L{0.8cm}L{1.3cm}L{7.1cm}L{5.5cm}@{}}
\caption{零密钥十分钟复现步骤}\label{tab:reproduce}\\
\toprule
\textbf{步骤} & \textbf{预算} & \textbf{命令或动作} & \textbf{通过标准} \\
\midrule
\endfirsthead
\toprule
\textbf{步骤} & \textbf{预算} & \textbf{命令或动作} & \textbf{通过标准} \\
\midrule
\endhead
1 & 0.5 min & \file{git clone https://github.com/Td12138/yizhi-chengke.git}，进入仓库 & 根目录可见 \file{apps/}、\file{docs/}、\file{README.md} \\
2 & 0.2 min & Windows 执行 \file{copy .env.example .env}；其他系统执行 \file{cp .env.example .env} & 根目录 \file{.env} 存在，两个 Provider 均为 mock \\
3 & 2--3 min & 在 \file{apps/server/} 创建 Python 3.11+ 虚拟环境并安装 \file{requirements.txt}；测试另装 \file{requirements-dev.txt} & 依赖安装无错误；不需要模型账号 \\
4 & 0.3 min & 在后端目录执行 \file{python -m uvicorn app.main:app --host 127.0.0.1 --port 8710} & 控制台显示监听 127.0.0.1:8710 \\
5 & 2--3 min & 另开终端，在 \file{apps/web/} 执行 \file{npm install} 和 \file{npm run dev} & Vite 显示 \url{http://127.0.0.1:5173} \\
6 & 0.2 min & 请求 \file{http://127.0.0.1:8710/api/health} & \env{ok=true}、两个 Provider 为 mock、\env{ready=true} \\
7 & 2 min & 打开页面，选择示例板书，生成 5 题，答错一题，再提交批改并追问 & 七 Agent 有状态；题目带级别；出现讲解、步骤诊断和引用 \\
8 & 可选 & 仓库根目录执行 \file{python -m pytest apps/server/tests -q} & 当前基线为 403 项通过；本机约 68 s \\
9 & 可选 & \file{apps/web/} 执行 \file{npm test} 与 \file{npm run build} & 18 项通过，类型检查与生产构建无错误 \\
\bottomrule
\end{longtable}

复现时建议先确认页面顶栏显示 \env{VLM: mock} 和 \env{LLM: mock}。缺省首轮 mock 命题会在第 4 题产生一次真实验算驳回，随后进入第二轮。若完全没有驳回，可能是样例或版本发生变化，应查看事件和测试，而不是手动伪造提示。若真实 Provider 显示未就绪，检查 Key、端点和模型名；系统不会把配置错误静默改成 mock。

演示前重置使用 \file{python apps/server/scripts/reset_demo.py}。该命令缺省保留原图，清理运行数据。需要删除归档图片时显式添加 \env{--purge-images}。桌面版使用包内重置脚本，执行前先退出程序。自动化测试自带临时数据目录，不需要为了跑测试而重置开发库。

\section{接口与配置表}

\subsection{HTTP 接口清单}

后端共有 17 个接口。异步任务接口立即返回 \env{run_id}，实际结果经 SSE 交付。查询接口同步返回，卡片修订是局部写操作。所有路径均以本机服务为预期部署环境，不包含鉴权字段。

\begin{longtable}{@{}L{1.2cm}L{4.7cm}L{4.7cm}L{4.5cm}@{}}
\caption{HTTP 接口、输入输出与用途}\label{tab:api}\\
\toprule
\textbf{方法} & \textbf{路径} & \textbf{请求到响应} & \textbf{说明} \\
\midrule
\endfirsthead
\toprule
\textbf{方法} & \textbf{路径} & \textbf{请求到响应} & \textbf{说明} \\
\midrule
\endhead
GET & \file{/api/health} & 六个就绪与 Provider 字段 & 配置与启动检查 \\
POST & \file{/api/ingest} & \env{multipart:file} $\to$ \env{run_id,kind} & 图片或文档识别、萃取与归档 \\
GET & \file{/api/events/{run_id}} & \env{text/event-stream} & 实时事件、迟到补发与保活 \\
GET & \file{/api/runs/{run_id}/events} & \env{run_id,events} & 已落库事件回放，无记录 404 \\
GET & \file{/api/cards} & \env{cards} & 最近 100 张卡片 \\
PATCH & \file{/api/cards/{id}} & 四个可选字段 $\to$ 完整卡片 & 白名单局部更新并留痕 \\
GET & \file{/api/cards/{id}/revisions} & \env{revisions} & 新到旧，最多 50 条 \\
GET & \file{/api/sources} & \env{sources} & 上传来源、卡片数与文件地址 \\
GET & \file{/api/sources/{id}} & 来源、原始块和卡片 & 不返回本机绝对路径 \\
GET & \file{/api/sources/{id}/image} & 文件二进制 & 按后缀返回 MIME，缺失 404 \\
GET & \file{/api/questions} & \env{questions} & 最近 50 道题 \\
POST & \file{/api/generate} & 卡片 ID、需求、题量 $\to$ \env{run_id} & 题量裁剪 1--12，空卡片取最近 3 张 \\
POST & \file{/api/grade} & 题号、作答 $\to$ \env{run_id} & 判分写 attempt，答错触发讲解 \\
POST & \file{/api/review} & 题号、作答 $\to$ \env{run_id} & 分步批改，结构结果随 done 返回 \\
POST & \file{/api/ask} & \env{query} $\to$ \env{run_id} & 检索作答并带出处 \\
GET & \file{/api/graph} & \env{nodes,edges,tags} & 共享标签边与卡片计数 \\
GET & \file{/api/stats} & \env{top_asked} & 高频卡片前 10 \\
\bottomrule
\end{longtable}

健康检查的六个字段为 \env{ok}、\env{vlm_provider}、\env{llm_provider}、\env{vlm_ready}、\env{llm_ready}、\env{ready}。真实 Provider 的就绪判断只检查 Key、地址和模型名是否齐全，不会为了健康检查额外调用付费接口。因此 \env{ready=true} 表示配置完整且服务已启动，不保证上游当时一定可用。

\subsection{主要配置项}

配置从仓库根 \file{.env} 或桌面版显式指定文件读取。已存在的系统环境变量优先，显式值优先于 Provider 预设。整数配置解析失败时退回默认并警告，带 BOM 的 UTF-8 文件和常见行尾注释可处理。密钥只应在本机配置，以下表格不包含任何真实值。

\begin{longtable}{@{}L{5.0cm}L{2.7cm}L{7.2cm}@{}}
\caption{主要环境变量、缺省值与回退行为}\label{tab:config}\\
\toprule
\textbf{变量} & \textbf{缺省值} & \textbf{作用与边界} \\
\midrule
\endfirsthead
\toprule
\textbf{变量} & \textbf{缺省值} & \textbf{作用与边界} \\
\midrule
\endhead
\env{VLM_PROVIDER} & mock & 视觉 Provider，可选 mock、qwen、glm；真实模式需 Key \\
\env{LLM_PROVIDER} & mock & 文本 Provider，可选 mock、deepseek、qwen、glm \\
\env{VLM_BASE_URL/API_KEY/MODEL} & 按预设或空 & 视觉端点、密钥与模型；显式值优先，密钥无预设 \\
\env{LLM_BASE_URL/API_KEY/MODEL} & 按预设或空 & 文本端点、密钥与模型；DeepSeek 只补文本预设 \\
\env{VERIFY_MAX_RETRY} & 2 & 首轮之后最多两次重试，总轮次最多 3 \\
\env{VERIFY_STRUCTURE_CHECK} & 1 & 无账单时做结构检查；设 0 恢复直接放行并标 unchecked \\
\env{VERIFY_DEGRADE_FORMAT} & 1 & 有限格式错误转结构检查；安全类拒绝不受影响 \\
\env{VERIFY_FALLBACK_ON_EXHAUST} & 1 & 重试耗尽且无通过项时，结构合格题可标级交付；设 0 严格拒绝 \\
\env{JUDGE_SYMPY_EQUIV} & 1 & 字符串与数值失败后追加符号等价；设 0 恢复旧规则 \\
\env{EVENT_PERSIST} & 1 & 事件落库；设 0 只保留内存事件 \\
\env{EVENT_PERSIST_RUNS} & 200 & 保留最近运行的数量，超过后裁剪旧事件 \\
\env{RETRIEVAL_MODE} & legacy & \env{legacy} 为 2-gram，\env{hybrid} 为向量与 RRF 可选增强 \\
\env{EMBED_BASE_URL/MODEL/API_KEY} & 空 & OpenAI 兼容嵌入端点；仅 hybrid 使用，故障回退 \\
\env{EMBED_TIMEOUT_SECONDS} & 10 & 嵌入请求超时 \\
\env{EMBED_BATCH_SIZE} & 16 & 单批嵌入文本数 \\
\env{EMBED_MAX_CHARS} & 1000 & 单卡送嵌入文本上限 \\
\env{HYBRID_CORPUS_LIMIT} & 500 & 向量路最多参与的卡片数 \\
\env{APP_HOST/APP_PORT} & 127.0.0.1 / 8710 & 后端监听；修改端口需同步开发代理 \\
\env{APP_DATA_DIR} & 仓库 data 或桌面用户目录 & 数据库与上传目录，测试会覆盖为临时路径 \\
\env{SERVE_WEB/WEB_DIST} & false / 前端 dist & 是否由后端托管已构建前端 \\
\env{WEB_CSP} & 1 & 托管前端时添加 CSP；设 0 仅供排障 \\
\env{YIZHI_ENV_FILE} & 空 & 桌面单文件版指向可执行文件旁配置 \\
\bottomrule
\end{longtable}

题量最小值、默认值和最大值分别为 1、5、12。图片与文档上限分别为 12 MB 和 20 MB，PDF 最多渲染前 3 页，单份文本最多 2 万字。检索固定返回前三条，向量余弦下限为 0.35，RRF 常数为 60。事件单 run 内存上限 500 条，保活间隔 15 秒。上述常量用于限制上下文、存储与资源消耗，不应被理解为已经经过大规模压测的最佳参数。

\section{测试矩阵}

本矩阵抽取主路径和高风险边界，便于评审不阅读全部测试源码也能复核覆盖意图。实际参数化用例多于表内行数，403 是 pytest 展开后的当前总数。命令与样本均使用自写数据，不包含真实密钥或个人信息。

\begin{longtable}{@{}L{1.3cm}L{3.4cm}L{5.5cm}L{5.0cm}@{}}
\caption{核心功能与异常路径测试矩阵}\label{tab:test-matrix}\\
\toprule
\textbf{编号} & \textbf{对象} & \textbf{输入或操作} & \textbf{期望结果与层级} \\
\midrule
\endfirsthead
\toprule
\textbf{编号} & \textbf{对象} & \textbf{输入或操作} & \textbf{期望结果与层级} \\
\midrule
\endhead
T01 & 上传后缀 & 不支持扩展名；超限图片或文档 & 400 或 413，不留下不完整文件；接口测试 \\
T02 & 正常图片 & PNG/JPG 等允许格式，12 MB 内 & 返回 \env{run_id} 与 image 类型；接口与 E2E \\
T03 & PDF 文本层 & 正常文本与少量空白尾页 & 文本层直读，不调用视觉模型；文档测试 \\
T04 & PDF 乱码 & 文本层不可读，真实 VLM 分支 & 前 3 页按 180 DPI 渲染后识别；MockTransport \\
T05 & 损坏文档 & 加密、零字节 PDF，损坏 DOCX & 即使 mock 也返回可读错误，不生成假卡片；单元测试 \\
T06 & DOCX/TXT/MD & 段落、表格、标题、不同编码 & 保持文档顺序与块类型；文档测试 \\
T07 & 原图归档 & 中文标题、非法字符、重名、无标题 & 清理名称、追加序号，无标题保留随机名；存储测试 \\
T08 & 卡片局部更新 & 改一个字段、提交同值、空 PATCH & 只改提交字段；同值不留历史；空修改幂等；接口测试 \\
T09 & 修订权限与历史 & 改来源字段、查不存在卡片、连续多次修改 & 来源字段不生效；404；历史新到旧且只记真实变化 \\
T10 & 命题数量 & 请求 0、5、20 或模型多返回 & 请求裁剪到 1--12，多返回部分截断；接口与单元测试 \\
T11 & 首轮驳回 & mock 第 1 轮第 4 题期望值为正 & 实际值 -1，驳回并进入第 2 轮；流水线测试 \\
T12 & 第二轮交付 & mock 第 2 轮全部账单自洽 & 命题共两次，五题交付，无运行错误；流水线测试 \\
T13 & 空结果与单题异常 & 模型返回空列表或某题字段类型损坏 & 不崩溃、不无限循环；异常限制在单题或明确运行错误 \\
T14 & 三级验算 & 有 verify、无 verify、关闭结构检查 & 分别得到 symbolic、structural、unchecked，界面与入库一致 \\
T15 & 结构检查 & 空题干、空答案、未知题型、坏选项、合法短题干 & 坏结构驳回，“解方程”等合法短题干不误伤 \\
T16 & 数值容差与期望 & 浮点误差、期望偏移、非数值或溢出期望 & 容差内通过，不自洽驳回，类型异常不拖垮整批 \\
T17 & 数学注入 & 动态导入、文件写入、eval/exec、属性链 & 全部拒绝，独立工作目录零新增文件；安全测试 \\
T18 & 代入与幂结构 & 表达式代入、布尔、容器、指数塔、嵌套幂 & 非数值拒绝；危险结构在求值前快速拒绝 \\
T19 & 正常数学兼容 & 多项式、sqrt、三角、log、pi、小指数 & 保持正确计算结果，防止白名单过严；单元测试 \\
T20 & 判分等价 & 方程交换、展开与因式、分数与小数、上标 & 等价答案判对；快速规则失败后走符号路径 \\
T21 & 判分反例 & \env{x=2} 对 \env{y=2}，\env{x/2=2} 对 2 & 均判错，不采用等号右端字面短路；单元测试 \\
T22 & 判分资源限制 & 高次幂、垃圾表达式、非法 LaTeX & 快速退回旧规则，不阻塞事件循环，不放宽结果 \\
T23 & 批改结构 & 完整、部分和错误作答；模型缺字段或错类型 & 分数、结论、步骤、首个错误完整且类型收敛 \\
T24 & 方法归因 & 字符串 false、空方法、另一种解法 & 布尔正确归一；空方法不输出；mock 不伪造归因 \\
T25 & 评语分片与回退 & 短尾片、中流异常、模型坏 JSON & 拼接不丢字不重叠；无“半句加全文”；事件标回退 \\
T26 & 检索基线 & “为什么这里要配方”与无关天气问题 & 前者命中正确卡片，后者返回空，不硬凑出处 \\
T27 & RRF 与缓存 & 两路排名、换模型、维度或有限数异常 & 共识卡片优先；换模型重嵌；坏向量拒绝 \\
T28 & 检索故障 & 无端点、超时、缓存数据库异常 & 自动回退 legacy，追问继续，事件注明路径 \\
T29 & 事件生命周期 & 迟到订阅、异常、容量超限、结束关闭 & 补发历史；终止事件保留；所有路径最终关闭 \\
T30 & 事件落库回放 & 写库失败、重启、旧 run 裁剪、未知 run & 实时流不受写库影响；重启可回放；未知返回 404 \\
T31 & 讲解参数 & 正常参数、NaN、字符串、反向区间 & 合法参数作图；非法参数只保留文字讲解 \\
T32 & 谱系与计数 & 泛化标签、无有效边、同批重复卡片、计数失败 & 过滤泛化边；同源链兜底；每卡只计一次；主任务继续 \\
T33 & 配置兼容 & BOM、行尾注释、坏整数、未知 Provider、显式覆盖 & 可解析或回默认；未知 Provider 不就绪；显式值优先 \\
T34 & 前端纯函数 & 上传闸门、题数、流式字段、批改清洗、回放合成 & 18 项通过，无浏览器和外部服务依赖 \\
T35 & 七条 E2E & 上传、出题、讲解、追问、批改、修订、回放 & 源码与打包版各 7/7；端口被占退出 2，不污染开发库 \\
\bottomrule
\end{longtable}

矩阵没有覆盖的项目同样需要记录：没有组件级视觉回归，没有多用户权限测试，没有高并发和长期稳定性测试，没有真实学生样本学习效果实验，没有所有 Provider 的真实联调。新增能力时应先在本矩阵加入可判定行，再实现功能。模型生成的候选用例必须经过人工复核，尤其要同时给出正例和反例。

\section{第三方依赖与许可}

下表列主要直接依赖和交付组件。传递依赖的完整许可文本以仓库和发布包中的 \file{THIRD_PARTY_NOTICES.md}、许可证文件及包管理锁文件为准。模型服务不是随源码再分发的软件组件，但其使用受各自服务条款约束。

\begin{longtable}{@{}L{3.5cm}L{3.0cm}L{4.2cm}L{4.0cm}@{}}
\caption{主要第三方组件、许可与用途}\label{tab:dependencies}\\
\toprule
\textbf{组件} & \textbf{许可} & \textbf{用途} & \textbf{交付边界} \\
\midrule
\endfirsthead
\toprule
\textbf{组件} & \textbf{许可} & \textbf{用途} & \textbf{交付边界} \\
\midrule
\endhead
FastAPI & MIT & HTTP 接口、参数校验与 ASGI 应用 & 后端直接依赖 \\
Uvicorn & BSD-3-Clause & ASGI 服务运行 & 单 worker 为当前运行方式 \\
python-multipart & Apache-2.0 & multipart 上传解析 & 上传接口直接依赖 \\
httpx & BSD-3-Clause & OpenAI 兼容模型与嵌入请求 & 支持 MockTransport 测试 \\
sympy & BSD-3-Clause & 受限表达式、代入验算与等价判分 & 不直接接受裸 \env{sympify} 输入 \\
pypdf & BSD-3-Clause & PDF 辅助处理与验证 & 后端直接依赖 \\
PyMuPDF & AGPL-3.0 或商业许可 & PDF 文本提取与页面渲染 & 桌面二进制包含，需保留对应义务 \\
python-docx & MIT & DOCX 段落与表格读取 & 后端直接依赖 \\
React & MIT & 单页界面与组件状态 & 前端运行时依赖 \\
TypeScript & Apache-2.0 & 类型检查与源码开发 & 构建工具，不是后端运行时 \\
Vite & MIT & 前端开发服务器与生产构建 & dist 不直接入库 \\
KaTeX & MIT & 题干与卡片公式渲染 & 错误时退回原文显示 \\
Electron & MIT & Windows 桌面壳、进程与配置定位 & 桌面分发组件 \\
PyInstaller & GPL-2.0+ 与打包例外 & 后端及 Python 运行时收集 & 按自身例外与 NOTICE 交付 \\
Playwright & Apache-2.0 & Chromium 端到端冒烟 & 开发与测试依赖，不是业务运行时 \\
SQLite & Public Domain & 本地结构化存储与 WAL & 随 Python 或运行环境使用 \\
Python & PSF License & 后端运行时 & 桌面包内置运行时并附许可 \\
Node.js & MIT & 前端构建与测试运行 & Web 开发环境需要，桌面用户无需另装 \\
\bottomrule
\end{longtable}

源码 MIT 许可只覆盖本项目可授权部分，不能覆盖第三方组件自身条款。特别是 PyMuPDF 的 AGPL 或商业许可选择需要在二进制分发时持续处理。发布包必须保留第三方 NOTICE、许可文本和对应源码获取说明。若后续更换 PDF 组件，应同步更新依赖锁、NOTICE、打包内容和本表，而不是只删除一行依赖。

\section{术语表}

\begin{longtable}{@{}L{3.4cm}L{11.5cm}@{}}
\caption{本文术语与项目内含义}\label{tab:glossary}\\
\toprule
\textbf{术语} & \textbf{定义} \\
\midrule
\endfirsthead
\toprule
\textbf{术语} & \textbf{定义} \\
\midrule
\endhead
Agent & 具有单一职责的异步处理单元。本文七个 Agent 不直接互相调用，由编排器传递状态 \\
Provider & 视觉或语言模型的可替换提供方。mock 与真实 API 实现同一调用边界 \\
mock & 内置自写样例 Provider。经过真实编排和规则，但视觉部分不读取图片语义 \\
Run / run\_id & 一次用户动作对应的运行及唯一编号，用于订阅、事件落库、恢复和排障 \\
RunContext & 单次运行的数据上下文，卡片、题目、作答和中间结果经它在 Agent 间传递 \\
SSE & Server-Sent Events。服务端向浏览器单向推送事件的 HTTP 流式协议 \\
delta & 增量事件，可携带逐题、逐卡或评语文本分片，不等同于最终结果 \\
rejected & 质检业务驳回状态，表示题目未通过当前检查，不表示系统异常 \\
verify 账单 & 命题模型提交的表达式、数值代入与期望值，用于本地独立核算 \\
symbolic & 符号验算级别，表示受限表达式已在给定代入下完成数值核算 \\
structural & 结构检查级别，只表示题干、答案、题型与选项形状符合要求 \\
unchecked & 未验算级别，保留关闭结构检查后的兼容行为，不形成内容质量证明 \\
safe\_math & 项目的受限数学解析入口，包含字符、标识符、代入值和幂结构白名单 \\
RRF & Reciprocal Rank Fusion，倒数排名融合。按两路检索名次合成分数，不直接混合异质量纲原始分 \\
2-gram & 中文相邻双字切片。项目用它避免依赖空白分词，并配合字段权重和分数下限 \\
BGE-M3 & 可选混合检索评测使用的现成嵌入模型，不是本作品训练的模型 \\
card\_revision & 卡片字段修订记录，保存旧值、新值和时间，与卡片更新在同一事务写入 \\
agent\_event & 落库运行事件，支持内存运行回收或服务重启后的结局恢复 \\
json-fallback & 批改模型或 JSON 路径不可用时回退本地规则的事件标记 \\
CSP & Content Security Policy。桌面壳和后端静态托管用于限制页面资源加载策略 \\
KaTeX & 前端公式渲染组件。公式语法失败时页面保留原文，不阻断整个界面 \\
WAL & SQLite 的预写日志模式，改善读写并发，但不等于支持任意规模服务 \\
结构降级 & 在限定条件下从符号验算退回结构检查，必须修改级别、记录原因并发出提示 \\
top-3 命中率 & 正确卡片是否出现在检索结果前三名的样本比例，不是答案文本准确率 \\
\bottomrule
\end{longtable}

\vspace{0.8em}
\begin{keybox}[文档结束与复核入口]
作品源码、复现命令、测试和公开材料统一从 \repo 获取。本文数字按 v0.1.6 现状整理。若代码后续变化，应同时更新版本、测试结果、配置表、接口表和边界说明，避免只修改功能清单。
\end{keybox}
